\documentclass[../DoAn.tex]{subfiles}
\begin{document}
\section{Đặt vấn đề}
\label{section:1.1}
Trong những năm gần đây, thương mại điện tử và các kênh bán hàng trực tuyến ngày càng trở thành phương thức tiếp cận khách hàng quan trọng đối với nhiều cửa hàng bán lẻ. Đối với lĩnh vực nội thất, quá trình mua sắm trực tuyến không chỉ dừng lại ở việc xem hình ảnh và giá bán, mà còn đòi hỏi người bán phải tư vấn dựa trên nhiều yếu tố như kích thước sản phẩm, chất liệu, phong cách thiết kế, không gian sử dụng và ngân sách của khách hàng. Vì vậy, hoạt động tư vấn đóng vai trò quan trọng trong việc giúp khách hàng hiểu rõ sản phẩm, so sánh các lựa chọn và đưa ra quyết định mua hàng phù hợp.

Tuy nhiên, đối với các cửa hàng nội thất nhỏ và vừa, việc duy trì một đội ngũ tư vấn có khả năng phản hồi nhanh, nhất quán và đủ hiểu biết về toàn bộ danh mục sản phẩm là một thách thức. Số lượng sản phẩm có thể thay đổi theo thời gian, thông tin sản phẩm thường bao gồm nhiều thuộc tính khác nhau, trong khi nhu cầu của khách hàng lại không cố định ngay từ đầu. Trong một phiên tư vấn, khách hàng có thể thay đổi ngân sách, điều chỉnh mục đích sử dụng, yêu cầu sản phẩm theo phong cách khác hoặc từ chối các gợi ý trước đó để tìm phương án phù hợp hơn. Điều này khiến quá trình tư vấn bán hàng trở thành một chuỗi trao đổi nhiều bước, cần khả năng ghi nhận ngữ cảnh và phản hồi linh hoạt.

Bên cạnh đó, khách hàng khi mua sản phẩm nội thất thường có nhu cầu so sánh giữa nhiều sản phẩm trước khi quyết định. Việc so sánh này không chỉ dựa trên giá, mà còn liên quan đến chất liệu, kích thước, kiểu dáng và mức độ phù hợp với không gian sử dụng. Nếu thông tin sản phẩm không được tổ chức tốt hoặc không được khai thác hiệu quả, khách hàng có thể mất nhiều thời gian để tìm kiếm, trong khi cửa hàng cũng khó tận dụng dữ liệu sản phẩm để nâng cao chất lượng tư vấn. Ngoài ra, nhu cầu tham chiếu mặt bằng giá thị trường cũng xuất hiện trong quá trình mua sắm, đặc biệt khi khách hàng muốn đánh giá một sản phẩm có phù hợp với ngân sách và giá trị sử dụng hay không.

Từ thực tế trên, có thể thấy bài toán hỗ trợ tư vấn bán hàng, so sánh sản phẩm và tham chiếu giá trong lĩnh vực nội thất có ý nghĩa thiết thực đối với cả người mua và cửa hàng. Nếu bài toán này được giải quyết hiệu quả, khách hàng có thể tiếp cận thông tin sản phẩm nhanh hơn, nhận được tư vấn phù hợp hơn với nhu cầu cá nhân và giảm thời gian ra quyết định. Đồng thời, các cửa hàng nội thất nhỏ và vừa có thể cải thiện chất lượng phục vụ, quản lý quá trình tư vấn nhất quán hơn và khai thác tốt hơn dữ liệu sản phẩm sẵn có. Đây là cơ sở thực tiễn để đồ án lựa chọn và triển khai đề tài liên quan đến hệ thống hỗ trợ tư vấn bán hàng trong lĩnh vực nội thất.

\section{Mục tiêu và phạm vi đề tài}
\label{section:1.2}
Hiện nay, các cửa hàng bán lẻ trực tuyến có thể sử dụng nhiều công cụ chatbot để hỗ trợ chăm sóc khách hàng và tư vấn bán hàng. Một nhóm phổ biến là các chatbot kịch bản như ManyChat hoặc Tidio. Các công cụ này phù hợp với các luồng hỏi đáp cố định, ví dụ như chào hỏi, thu thập thông tin liên hệ, trả lời một số câu hỏi thường gặp hoặc chuyển tiếp cho nhân viên. Tuy nhiên, điểm yếu của nhóm công cụ này là phụ thuộc nhiều vào các flow được cấu hình trước. Khi khách hàng thay đổi nhu cầu, đổi ngân sách, từ chối sản phẩm đã gợi ý hoặc đặt câu hỏi theo cách không nằm trong kịch bản, chatbot khó duy trì mạch tư vấn và thường phải quay lại các nhánh hội thoại đơn giản.

Một nhóm giải pháp khác là các chatbot có khả năng truy xuất dữ liệu từ knowledge base, ví dụ như Botpress hoặc các nền tảng chatbot có hỗ trợ tài liệu nội bộ. Nhóm này phù hợp hơn với các bài toán hỏi đáp dựa trên dữ liệu có sẵn, vì chatbot có thể tìm kiếm thông tin trong tài liệu hoặc dữ liệu sản phẩm trước khi trả lời. Tuy nhiên, việc xây dựng, cập nhật và quản lý knowledge base thường yêu cầu người dùng có hiểu biết nhất định về cấu trúc dữ liệu, tài liệu, pipeline xử lý và cách kiểm soát chất lượng câu trả lời. Đối với các cửa hàng nội thất nhỏ và vừa, đây là một rào cản đáng kể, vì họ thường không có đội ngũ kỹ thuật riêng để vận hành một hệ thống AI hoàn chỉnh.

Bên cạnh đó, các chatbot AI thương mại như Intercom Fin có khả năng trả lời tự nhiên và tích hợp vào quy trình chăm sóc khách hàng của doanh nghiệp. Tuy nhiên, các nền tảng này thường gắn với hệ sinh thái riêng, có chi phí vận hành tăng theo mức sử dụng và không phải lúc nào cũng phù hợp với nhu cầu xây dựng một hệ thống có dữ liệu riêng cho từng cửa hàng. Trong bối cảnh đồ án, nhóm giải pháp này cũng chưa trực tiếp giải quyết bài toán một nền tảng dùng chung cho nhiều cửa hàng nội thất, trong đó mỗi cửa hàng cần có chatbot, dữ liệu sản phẩm, cấu hình và lịch sử hội thoại riêng.

Đối với lĩnh vực nội thất, nhu cầu tư vấn bán hàng có một số đặc điểm cụ thể. Sản phẩm thường có nhiều thuộc tính ảnh hưởng trực tiếp đến quyết định mua, gồm tên sản phẩm, mô tả, giá, chất liệu, kích thước, loại sản phẩm, phong cách, cửa hàng cung cấp và liên kết sản phẩm nếu có. Người dùng cũng không phải lúc nào xác định rõ nhu cầu ngay từ đầu. Trong một phiên tư vấn, họ có thể bắt đầu bằng một yêu cầu chung, sau đó điều chỉnh ngân sách, thay đổi mục đích sử dụng, yêu cầu gợi ý lại, so sánh giữa nhiều lựa chọn hoặc hỏi thêm về mặt bằng giá. Vì vậy, bài toán không chỉ là trả lời một câu hỏi đơn lẻ, mà là hỗ trợ một quá trình tư vấn nhiều lượt có ngữ cảnh.

Từ các phân tích trên, có thể khái quát ba hạn chế chính của các giải pháp hiện tại đối với bài toán trong đồ án. Thứ nhất, chatbot kịch bản chưa phù hợp với các tình huống tư vấn sản phẩm có nhiều thay đổi trong hội thoại. Thứ hai, chatbot có knowledge base hoặc RAG thường đòi hỏi khả năng kỹ thuật khi xây dựng và cập nhật dữ liệu, gây khó khăn cho cửa hàng nhỏ và vừa. Thứ ba, các nền tảng chatbot AI thương mại có chi phí và ràng buộc tích hợp, trong khi bài toán cần một hệ thống có thể phục vụ nhiều cửa hàng, tách biệt dữ liệu theo từng cửa hàng và vẫn hỗ trợ so sánh, tham chiếu trên dữ liệu tổng hợp.

Trên cơ sở đó, mục tiêu của đề tài là xây dựng một hệ thống chatbot hỗ trợ tư vấn bán hàng, so sánh sản phẩm và tham chiếu giá cho lĩnh vực nội thất. Hệ thống hướng đến hai nhóm người dùng chính. Nhóm thứ nhất là các cửa hàng nội thất nhỏ và vừa, có nhu cầu sử dụng chatbot trên dữ liệu sản phẩm riêng nhưng không muốn tự xây dựng một hệ thống AI độc lập. Nhóm thứ hai là người mua hàng, có nhu cầu được tư vấn theo ngôn ngữ tự nhiên, được gợi ý sản phẩm phù hợp, so sánh các sản phẩm theo tiêu chí rõ ràng và tham khảo mặt bằng giá trước khi ra quyết định.

Về mục tiêu chức năng, hệ thống cần cung cấp các thành phần phục vụ quản lý tenant, chatbot, người dùng, hội thoại, knowledge base và yêu cầu mua hàng. Đối với cửa hàng, hệ thống cần cho phép quản lý dữ liệu và cấu hình chatbot, theo dõi hội thoại, kiểm tra trạng thái knowledge base và tiếp nhận các yêu cầu mua hàng được tạo trong quá trình tư vấn. Đối với người dùng cuối, hệ thống cần hỗ trợ hội thoại nhiều lượt, ghi nhận các thông tin đã trao đổi, điều chỉnh gợi ý khi yêu cầu thay đổi và tạo yêu cầu mua hàng khi đã có đủ thông tin cần thiết.

Phạm vi hoạt động của hệ thống gồm ba chế độ chính. Chế độ thứ nhất là tư vấn theo cửa hàng, trong đó chatbot sử dụng dữ liệu riêng của từng cửa hàng để tư vấn sản phẩm, duy trì mạch hội thoại và hỗ trợ tạo yêu cầu mua hàng. Đây là chế độ phục vụ trực tiếp cho hoạt động bán hàng của từng cửa hàng và có thể được triển khai qua các kênh như web chat, Messenger hoặc Telegram. Chế độ thứ hai là tư vấn và so sánh chung trên website của hệ thống, trong đó chatbot sử dụng dữ liệu tổng hợp từ các cửa hàng tham gia để trả về các sản phẩm phù hợp và trình bày các tiêu chí so sánh như giá, chất liệu, kích thước, loại sản phẩm hoặc phong cách. Chế độ này chỉ phục vụ nhu cầu tham khảo và so sánh, không tạo lead hoặc yêu cầu mua hàng cho một cửa hàng cụ thể.

Chế độ thứ ba là tham chiếu thị trường, trong đó hệ thống sử dụng dữ liệu sản phẩm tổng hợp để cung cấp thông tin về khoảng giá, mức giá phổ biến và các trường hợp có giá lệch đáng kể so với nhóm sản phẩm tương ứng. Phạm vi của chế độ này không phải là định giá toàn bộ thị trường nội thất, mà là hỗ trợ người dùng tham khảo mặt bằng giá dựa trên dữ liệu hiện có trong hệ thống. Chức năng này phù hợp với các tình huống người dùng muốn biết một loại sản phẩm thường có giá trong khoảng nào hoặc một sản phẩm cụ thể đang cao hay thấp hơn đáng kể so với nhóm tương tự.

Trong phạm vi đồ án, hệ thống cần được xây dựng dưới dạng phần mềm có backend nghiệp vụ, AI service xử lý hội thoại và truy xuất tri thức, giao diện quản trị cho cửa hàng, giao diện quản trị hệ thống và giao diện chat cho người dùng cuối. Dữ liệu demo cần bao gồm tối thiểu hai tenant hoặc cửa hàng mẫu, mỗi tenant có dữ liệu sản phẩm hoặc tài liệu sản phẩm riêng. Hệ thống cũng cần có dữ liệu tổng hợp để phục vụ chức năng so sánh chung và tham chiếu giá.

Việc đánh giá đề tài tập trung vào các tiêu chí cụ thể đã được xác định trong phiếu giao nhiệm vụ. Đối với chế độ tư vấn theo cửa hàng, hệ thống cần được đánh giá qua mức độ phù hợp của câu trả lời với ngữ cảnh và dữ liệu sản phẩm, khả năng xử lý các tình huống hội thoại đa lượt như đổi ngân sách, đổi nhu cầu, từ chối đề xuất, và khả năng tạo yêu cầu mua hàng khi đủ thông tin. Đối với chế độ so sánh chung, hệ thống cần được đánh giá qua khả năng trả về các sản phẩm phù hợp và trình bày được các tiêu chí so sánh chính. Đối với chế độ tham chiếu thị trường, hệ thống cần được đánh giá qua khả năng tính khoảng giá cho các nhóm sản phẩm có đủ dữ liệu và phát hiện các trường hợp giá bất thường. Ngoài ra, đồ án cũng thực hiện so sánh với một số mô hình ngôn ngữ lớn hoặc chatbot có công dụng tương tự để có cơ sở nhận xét về chất lượng kết quả đạt được.

Như vậy, phạm vi của đề tài không đặt mục tiêu xây dựng một nền tảng thương mại điện tử hoàn chỉnh hay bao phủ toàn bộ thị trường nội thất. Trọng tâm của đề tài là xây dựng và đánh giá một hệ thống chatbot phục vụ tư vấn bán hàng nội thất trong môi trường nhiều cửa hàng, có dữ liệu riêng theo từng cửa hàng, có khả năng hỗ trợ hội thoại nhiều lượt, so sánh sản phẩm và tham chiếu giá dựa trên dữ liệu hiện có.

\section{Định hướng giải pháp}
\label{section:1.3}
Từ các mục tiêu và phạm vi đã xác định, đồ án lựa chọn định hướng xây dựng hệ thống chatbot tư vấn bán hàng dựa trên mô hình đa tenant kết hợp với kỹ thuật Retrieval-Augmented Generation. Định hướng này phù hợp với bài toán vì mỗi cửa hàng nội thất cần có dữ liệu, chatbot và cấu hình riêng, trong khi hệ thống vẫn cần vận hành trên một nền tảng thống nhất. Bên cạnh đó, RAG cho phép chatbot sinh câu trả lời dựa trên dữ liệu được truy xuất từ knowledge base, qua đó hạn chế việc trả lời thiếu căn cứ và giúp nội dung tư vấn bám sát dữ liệu sản phẩm hiện có.

Theo định hướng trên, giải pháp của đồ án là xây dựng một hệ thống gồm hai nhóm thành phần chính: phần nghiệp vụ quản lý cửa hàng và phần AI phục vụ hội thoại tư vấn. Phần nghiệp vụ chịu trách nhiệm quản lý tenant, chatbot, người dùng, hội thoại, knowledge base và yêu cầu mua hàng. Phần AI service xử lý hội thoại, truy xuất thông tin từ knowledge base, đưa ngữ cảnh phù hợp vào prompt và sinh câu trả lời tư vấn cho người dùng. Hệ thống được thiết kế để hỗ trợ ba chế độ hoạt động đã xác định trong phạm vi đề tài, gồm tư vấn theo cửa hàng, tư vấn và so sánh chung, và tham chiếu mặt bằng giá.

Đối với bài toán truy xuất thông tin, đồ án định hướng sử dụng các phương pháp retrieval gồm TF-IDF hoặc keyword retrieval làm baseline, vector retrieval để khai thác độ tương đồng ngữ nghĩa, và hybrid retrieval để so sánh kết quả giữa các cách tiếp cận. Các phương pháp này được lựa chọn vì phù hợp với dữ liệu sản phẩm nội thất, trong đó thông tin có thể vừa mang tính thuộc tính cụ thể như giá, chất liệu, kích thước, vừa mang tính mô tả như phong cách hoặc mục đích sử dụng. Kết quả truy xuất được dùng làm cơ sở để chatbot đưa ra câu trả lời, gợi ý sản phẩm hoặc hỗ trợ so sánh sản phẩm.

Về công nghệ triển khai, đồ án định hướng xây dựng backend nghiệp vụ bằng Java và Spring Boot, AI service bằng Python và FastAPI, cơ sở dữ liệu bằng PostgreSQL, đồng thời đóng gói các thành phần bằng Docker và Docker Compose. Các công nghệ này được lựa chọn nhằm tách biệt rõ phần nghiệp vụ và phần xử lý AI, giúp hệ thống dễ phát triển, kiểm thử và triển khai trong môi trường demo nhiều service. Đối với kênh tương tác, hệ thống hướng tới hỗ trợ web chat và tích hợp với các kênh phổ biến như Facebook Messenger và Telegram để phù hợp với cách các cửa hàng nhỏ và vừa thường giao tiếp với khách hàng.

Đóng góp chính của đồ án là thiết kế và xây dựng một hệ thống chatbot RAG đa tenant cho lĩnh vực tư vấn bán hàng nội thất. Hệ thống cho phép mỗi cửa hàng vận hành chatbot trên dữ liệu riêng, hỗ trợ hội thoại nhiều lượt, ghi nhận sự thay đổi trong nhu cầu của khách hàng, gợi ý và so sánh sản phẩm dựa trên dữ liệu, đồng thời tạo yêu cầu mua hàng khi có đủ thông tin. Ngoài chế độ tư vấn theo từng cửa hàng, hệ thống còn cung cấp chế độ tư vấn và so sánh chung trên dữ liệu tổng hợp, cũng như chế độ tham chiếu giá để hỗ trợ người dùng đánh giá mặt bằng giá của các nhóm sản phẩm.

Kết quả cuối cùng của đồ án hướng tới một sản phẩm phần mềm có thể chạy thử trong môi trường demo với dữ liệu mẫu của tối thiểu hai cửa hàng, knowledge base theo từng cửa hàng, giao diện quản trị, giao diện chat và các API phục vụ hội thoại, quản lý dữ liệu, so sánh sản phẩm và tham chiếu giá. Bên cạnh sản phẩm phần mềm, đồ án cũng xây dựng tập benchmark và kịch bản kiểm thử để đánh giá chất lượng truy xuất, chất lượng câu trả lời, khả năng xử lý hội thoại đa lượt, khả năng tạo yêu cầu mua hàng và khả năng hỗ trợ so sánh, tham chiếu giá.

\section{Bố cục đồ án}
\label{section:1.4}
Phần còn lại của báo cáo đồ án tốt nghiệp này được tổ chức như sau.

\textbf{Chương 2} trình bày quá trình khảo sát bài toán và phân tích nhu cầu của các nhóm người dùng trong hệ thống chatbot tư vấn bán hàng nội thất. Nội dung chương tập trung làm rõ bối cảnh sử dụng hệ thống, các đối tượng liên quan như người dùng cuối, cửa hàng và quản trị hệ thống, đồng thời phân tích các yêu cầu chức năng và phi chức năng cần thiết. Chương này cũng trình bày phạm vi dữ liệu được sử dụng trong hệ thống, các tình huống hội thoại tiêu biểu và những tiêu chí cần xem xét khi đánh giá chất lượng tư vấn, so sánh sản phẩm và tham chiếu giá.

\textbf{Chương 3} giới thiệu các công nghệ và cơ sở kỹ thuật được sử dụng để xây dựng hệ thống. Chương này trình bày tổng quan về các thành phần công nghệ chính phục vụ phát triển backend, AI service, cơ sở dữ liệu, giao diện người dùng và môi trường triển khai. Bên cạnh đó, chương cũng giới thiệu các kỹ thuật liên quan đến chatbot, knowledge base, truy xuất thông tin và Retrieval-Augmented Generation ở mức tổng quan, nhằm làm cơ sở cho việc thiết kế và triển khai hệ thống trong các chương tiếp theo.

\textbf{Chương 4} trình bày kết quả thực nghiệm và đánh giá hệ thống. Nội dung chương tập trung vào cách xây dựng dữ liệu mẫu, tập câu hỏi benchmark và các kịch bản kiểm thử phục vụ đánh giá. Trên cơ sở đó, chương phân tích kết quả của hệ thống theo các nhóm tiêu chí chính, gồm chất lượng truy xuất từ knowledge base, chất lượng câu trả lời tư vấn, khả năng xử lý hội thoại nhiều lượt, khả năng gợi ý và so sánh sản phẩm, khả năng tham chiếu giá và khả năng tạo yêu cầu mua hàng. Chương này cũng trình bày các nhận xét về điểm đạt được và những hạn chế còn tồn tại sau quá trình thử nghiệm.

\textbf{Chương 5} trình bày giải pháp xây dựng hệ thống và các đóng góp chính của đồ án. Chương này mô tả cách hệ thống được tổ chức để hỗ trợ nhiều cửa hàng trên cùng một nền tảng, trong đó mỗi cửa hàng có dữ liệu, chatbot và cấu hình riêng. Nội dung chương cũng trình bày các thành phần chính của hệ thống, gồm backend nghiệp vụ, AI service, giao diện quản trị, giao diện chat và các nhóm API phục vụ hội thoại, quản lý knowledge base, quản lý yêu cầu mua hàng, so sánh sản phẩm và tham chiếu giá. Đây là chương tập trung làm rõ sản phẩm phần mềm được xây dựng trong đồ án.

\textbf{Chương 6} tổng kết các nội dung đã thực hiện trong đồ án và đánh giá mức độ hoàn thành so với mục tiêu ban đầu. Chương này trình bày những kết quả chính đã đạt được, các hạn chế còn tồn tại trong phạm vi triển khai và đánh giá, đồng thời đề xuất một số hướng phát triển tiếp theo cho hệ thống. Các hướng phát triển có thể bao gồm mở rộng dữ liệu, cải thiện chất lượng truy xuất và câu trả lời, bổ sung kênh tích hợp, nâng cao khả năng quản trị knowledge base và hoàn thiện hệ thống để tiến gần hơn tới môi trường vận hành thực tế.

\end{document}